
% --------------------------------------------------------------
%
%  Modelo de artigo do 10º MCSul / X SEMENGO (FURG) para LaTeX
%
%  O modelo de artigo consiste dos seguintes arquivos:
%   - template.tex    > Este arquivo
%   - semengo.cls     > Definição da classe de documento
%   - logo.jpg        > Logotipo do evento
%   - references.bib  > Referências de exemplo
%  E do guia de usuário:
%   - semengo-doc.pdf
%
%  Este documento pode ser usado como base para escrever seu
%  trabalho. Ele está comentado para auxiliar o autor.
%
%  Bugs podem ser relatados por e-mail:
%    phelype.oleini

%@latex-project.org
%
% --------------------------------------------------------------

% Use a opção global da classe para mudar o idioma do
% documento. Use 'brazilian' (ou 'portuguese') para
% português (europeu), ou 'USenglish' (ou 'UKenglish')
% para inglês (britânico).
\documentclass[brazilian]{semengo}

% Insira o título com \title:
\title{Guia de Submissão para Autores}

% Insira os autores, separados por \and com \author.
% Em \author, use \addr{label} para inserir o endereço
% de cada autor, e \corr{email} para adicionar o
% endereço eletrônico do autor correspondente.
\author{ Nome do Autor \addr{endereço-1} \corr{endereco@eletronico.com}
  \and Nome dos Coautores \addr{endereço-2} }

% Defina os endereços usando \address. Cada endereço
% é definido com "label = endereço", separados por \and:
\address{ endereço-1 = { Instituição/endereço profissional 1 }
     \and endereço-2 = { Instituição/endereço profissional 2 } }

% Identifique a imagem do logotipo usando \mainlogo:
\mainlogo{logo2024.png}

% Escreva o resumo do trabalho no ambiente 'abstract'.
% certifique-se que '\begin{abstract}' seja o único
% comando em sua linha, assim como '\end{abstract}',
% ou coisas ruins podem acontecer :-)
\begin{abstract}
O resumo deverá compreender, de forma concisa, a temática abordada e
apresentar em sua estrutura os objetivos, bem como a metodologia
aplicada, e uma prévia dos resultados alcançados. Deve ser escrito em um
único parágrafo. Insira o resumo usando o ambiente \texttt{abstract}
(\verb|\begin{abstract}...\end{abstract}|) no preâmbulo do documento
(antes de \verb|\begin{document}|). O resumo deve ter entre 200 e 300
palavras. Logo após o resumo, de três a cinco palavras-chave devem
destacar de modo claro a temática do artigo. Insira as palavras-chave
usando \verb|\keywords{keyword, keyword, ...}|. O título, lista de
autores, resumo, palavras-chave, e notas de rodapé da página inicial não
devem avançar para a segunda página. O idioma do trabalho deve ser
especificado na lista de opções globais:
\verb|\documentclass[english]{semengo}|. Se o trabalho for escrito em
português use \texttt{portuguese} ou \texttt{brazilian}, e se for em
inglês use \texttt{UKenglish} ou \texttt{USenglish}, para os dialetos
europeus ou americanos, respectivamente. Outros idiomas também são
possíveis: verifique a disponibilidade no manual do pacote
\textsf{babel}.
\end{abstract}

% Insira as palavras-chave separadas por vírugula
% usando \keywords:
\keywords{Palavra-chave, palavra-chave,
  palavra-chave, palavra-chave, palavra-chave}

% Indique o banco de referências (.bib) com:
\addbibresource{references.bib}

% Alguns pacotes básicos
% (insira ou remova conforme necessário).

% amsmath para equações
\usepackage{amsmath}

% siunitx para unidades do SI (\SI{numero}{unidade})
\usepackage{siunitx}

% booktabs para tabelas (\toprule, \bottomrule, ...)
\usepackage{booktabs}

% subcaption para sub-figuras
\usepackage{subcaption}

% hypreref e bookmark para hyperlinks
\usepackage[hidelinks]{hyperref}
\usepackage{bookmark}
% hyperref é (via de regra) o último pacote a ser
% carregado, à exceção de 'cleveref', 'amsrefs',
% e poucos outros.
% (https://tex.stackexchange.com/q/1863/134574)

% Configuração do pacote siunitx:
% Declarando a unidade \point:
\DeclareSIUnit{\point}{pt}
\sisetup{
  , detect-all % detectar fontes
  , output-decimal-marker = { , } % vírgula como separador decimal
}
%

% Configuração do pacote amsmath:
% Declarando um operador matemático:
\DeclareMathOperator{\Var}{Var}

% Início do corpo do texto
\begin{document}

\section{Introdução}

A introdução é o texto de abertura do artigo. Ela deve mostrar a
relevância do problema científico que motivou a realização do trabalho
para área, apresentar a finalidade e os objetivos do trabalho
proposto, destacando o aspecto inédito do trabalho e sua contribuição
para área, bem como outras questões relevantes que explorem a temática
abordada. Portanto, segue uma estrutura que permita delinear o assunto
foco do trabalho, as referências que dão embasamento à discussão do
tema, a justificativa para a escolha do tema e da importância do
trabalho, o problema, os objetivos e o método utilizado.

Para subdividir o documento utilize seções, subseções e sub-subseções,
conforme necessário. A seguir é mostrado um exemplo de subseção,
e de sub-subseção. Não devem ser usados níveis inferiores de subdivisão
do documento. Deve-se evitar também títulos de seção consecutivos sem
ao menos um parágrafo de texto entre eles, como é feito mais adiante na
seção~\ref{sec:figs_and_tabs} e subseção~\ref{sec:figures}.

\subsection{Formatação do Texto}

O corpo do texto utiliza a fonte \emph{\TeX{} Gyre Termes} (clone da
fonte \emph{Times New Roman}) em tamanho \SI{12}{\point}, com
espaçamento de \num{1.2} entre linhas, e alinhamento justificado.

\subsubsection{Instruções Adicionais}

Os anais do 10º MCSUL e X SEMENGO serão publicados em formato
digital (PDF---Portable Document Format). Os artigos devem ser
formatados estritamente de acordo com as instruções aqui fornecidas.
O presente arquivo (\texttt{.tex}) pode ser utilizado como base para
o documento, ou como referência para outros programas de processamento
de texto, como L\kern-.1667em\lower.25em\hbox{Y}\kern-.125emX,
\emph{Microsoft Word}, ou \emph{LibreOffice Writer}, conforme
preferência do autor. De toda maneira, pode ser utilizado como guia de
formatação para usuários de outros softwares.

Artigos completos devem conter até~15 páginas, incluindo tabelas e
figuras. Os arquivos não devem ultrapassar 2\,MB.

\section{Formato do Texto}

Neste momento, apresentam-se com maiores detalhes a fundamentação
teórica do trabalho---detalhando o estado da arte sobre o problema---,
a metodologia, os resultados ou a discussão que couber
(contextualizada frente à literatura da área), o que pode variar de
acordo com a natureza de cada trabalho. Quando couber, as descrições
presentes na seção devem apresentar detalhamento suficiente para a
reprodução do trabalho por outros pesquisadores. Se os métodos forem
correntes e conhecidos podem ser descritos brevemente com uso de
referências, mas se forem originais, devem ser descritos
detalhadamente.

Quando cabível, a discussão dos resultados deve buscar explicações e
fazer comparações com a bibliografia da área de forma a dar ao
trabalho uma abrangência nacional ou global, não sendo apenas
confirmatório.

Uma estrutura que pode ser tomada como base é fazer o desenvolvimento do
trabalho com a revisão bibliográfica ocupando de \SI{10}{\%} a
\SI{20}{\%} da extensão do artigo; os materiais e métodos de
\SI{10}{\%} a \SI{20}{\%} e a discussão dos resultados ocupando de
\SI{40}{\%} a \SI{60}{\%}.

\section{Citações}

As citações e referências seguem o manual de estilo da APA\footnote
{\url{https://apastyle.apa.org/instructional-aids/reference-guide.pdf}}
(\emph{American Psychological Association}---APA,~\citeyear{APA2020}).
As citações são no formato autor-data, portanto citações narrativas são
no formato Autor~(ano), e citações entre parênteses são (Autor,~ano).

No caso de um único autor, o formato é Autor~(ano), e (Autor,~ano),
para citações narrativas e em parênteses, respectivamente, como em
\textcite{knuth:ct:a} ou \parencite{worman}. No caso de dois autores,
os sobrenomes são separados por \enquote{e}, ou por \enquote{and},
dependendo do idioma do corpo do texto, como \textcite{sigfridsson},
ou \parencite{brandt}. De forma similar, no caso de três autores, como
em \textcite{companion}, ou em \parencite{padhye}. Em citações com
mais de três autores, coloca-se o sobrenome do autor principal seguido
de \enquote{\emph{et al.}}, como em \textcite{murray} ou em
\parencite{yoon}.

Quando citando uma lista de trabalhos de um mesmo autor, deve-se citar
o sobrenome do autor apenas uma vez, seguido da lista de anos das
diferentes publicações. Por exemplo, ao invés de citar
\textcite{averroes/hercz} e \textcite{averroes/hannes}, cita-se
\textcite{averroes/hercz,averroes/hannes}. Da mesma se a citação for
entre parênteses \parencite{averroes/hercz,averroes/hannes}.

Ainda, quando o mesmo autor publicou dois trabalhos diferentes em um
mesmo ano, numera-se os trabalhos utilizando letras minúsculas, por
exemplo, \textcite{knuth:ct:b} e depois \textcite{knuth:ct:d}. Além
disso, a regra anterior se aplica nesse caso, mas usa-se apenas as
letras separadas por vírgula. Nesse caso, ao invés de citar
\textcite{knuth:ct:a}, \textcite{knuth:ct:b}, e \textcite{knuth:ct:d},
cita-se \textcite{knuth:ct:a,knuth:ct:b,knuth:ct:d}.

\section{Equações}

As equações são criadas usando os ambientes de equação padrão do
pacote \textsf{amsmath}: \texttt{equation} para equações simples,
\texttt{gather} para conjuntos de equações, e \texttt{align} também
para conjuntos equações, com alinhamento horizontal, entre outros.
Para omitir a numeração em um ambiente \texttt{env}, use a versão com
\verb|*|, \texttt{env*}. A equação~\eqref{eq:exemplo} é dada como
exemplo:
\begin{equation}
  x = \sum_{i=0}^z 2^iQ
  \label{eq:exemplo}
\end{equation}
onde $z$ e $i$ são variáveis inteiras, e $Q$ é uma constante, e essa
frase demonstrou como mencionar variáveis no texto. Equações devem ser
referenciadas no texto usando o comando~\verb|\eqref|, como em:
equação~\eqref{eq:exemplo}, ou entre parênteses como
(equação~\eqref{eq:exemplo}). A critério do autor, pacotes
especializados de referências cruzadas podem ser usados, como
o~\textsf{cleveref}.

Equações menores podem ser inseridas na linha com o texto, conforme o
exemplo a seguir. Sendo $X_1,X_2,\dots,X_n$ uma sequência de variáveis
independentes e identicamente distribuídas aleatoriamente com
$\operatorname{E}[X_i]=\mu$ e $\Var[X_i]=\sigma^2<\infty$,
\begin{equation}
 S_n = \frac{X_1+X_2+\dots+X_n}{n} = \frac{1}{n}\sum_{i=1}^n X_i
 \label{eq:outro-exemplo}
\end{equation}
denota seu significado. Então, quando $n$ tende ao infinito, as
variáveis aleatórias $\sqrt{n}(S_n-\mu)$ convergem em uma distribuição
normal $\mathcal{N}(0,\sigma^2)$. Isto é apresentado na
equação~\eqref{eq:outro-exemplo}.

\section{Figuras e Tabelas}
\label{sec:figs_and_tabs}

\subsection{Figuras}
\label{sec:figures}

As figuras devem ser inseridas no texto sempre após o ponto em que são
referenciadas pela primeira vez. Para inserir figuras numeradas utilize
o ambiente \texttt{figure}. Um exemplo comum seria como:
\begin{verbatim}
\begin{figure}
  \centering
  \includegraphics[width=0.789\textwidth]{example-image}
  \caption{Legenda da figura}
  \label{rótulo-para-referência}
\end{figure}
\end{verbatim}
No exemplo, o comando \verb|\centering| faz com que a figura seja
horizontalmente centralizada na área da página (caso contrário, seria
alinhada à margem esquerda). O comando \verb|\includegraphics| é o que
de fato inclui a figura no documento. A opção
\verb|width=0.789\textwidth| faz com que a figura tenha largura igual
a $0.789\times\verb|\textwidth|$, onde \verb|\textwidth| é a largura
da linha de texto. Usando essa opção, a altura da figura será alterada
automaticamente de modo a manter a proporção altura/largura.

Assim como foi dada a opção \verb|width|, pode-se usar a opção
\verb|height| para definir a altura da figura. No entanto, se usar
\emph{ambas} as opções \verb|width| e \verb|height|, a proporção da
figura não necessariamente se manterá.

Ainda no exemplo acima, o comando \verb|\caption| insere a legenda
abaixo da figura, e o comando \verb|\label{<rótulo>}| atribui um
\verb|<rótulo>| à legenda da figura para referência posterior
(usando \verb|\ref{<rótulo>}|). Note que o elemento numerado é a
\emph{legenda} da figura (\verb|\caption|), e não a figura em si,
portanto uma figura sem legenda não pode ser referenciada. Além disso,
em uma figura a \verb|\label| deve \emph{sempre} vir após
\verb|\caption|.

Finalmente, evite usar o argumento opcional do ambiente
\texttt{figure}, pois esta classe de documento já define um padrão
adequado. Use o argumento opcional apenas se necessário e se tiver
certeza do que está fazendo. Além disso, \emph{não use} o argumento
opcional \verb|[H]|.

Figuras devem ser referenciadas no texto como figura~\ref{fig:exemplo}
ou entre parênteses, como (figura~\ref{fig:exemplo}).
A figura~\ref{fig:exemplo} é um exemplo simples de como inserir uma
figura.

\begin{figure}
  \centering
  \includegraphics[width=0.5\textwidth]{example-image}
  \caption{Diagrama do aparato experimental.}
  \label{fig:exemplo}
\end{figure}

Para inserir sub-figuras recomenda-se o pacote \textsf{subcaption},
que define o ambiente \texttt{subfigure}. A seguir um exemplo de como
inserir uma figura com várias sub-figuras no texto. O código fonte
deste documento contém comentários que podem auxiliar.

\begin{figure} % sempre use indentação no código :-)
  \centering % para centralizar
  % O TeX vê as subfiguras da mesma forma que caracteres
  % (dadas as devidas proporções). Fica mais fácil se você
  % pensar assim também.
  % Primeira sub-figura (topo, esquerda)
  \begin{subfigure}[t]{0.3\textwidth} % <- Essa subfigura terá essa largura
    %  Use \textwidth aqui V--------V para a figura ter a mesma largura
    % do ambiente subfigure (0.3\textwidth, neste caso).
    \includegraphics[width=\textwidth]{example-image-a}
    \caption{Legenda da sub-figura 1\strut}
    \label{label-da-subfigura-1}
  \end{subfigure}% <- esse '%' deve ficar aqui para eliminar espaço
  % entre sub-figuras (veja: https://tex.stackexchange.com/a/407620/134574)
  \begin{subfigure}[t]{0.7\textwidth} % <- 0.7 vezes a largura do texto
    \includegraphics[height=3.35cm,width=\textwidth]{example-image-b}
    \caption{Legenda da sub-figura 2\strut}
    \label{label-da-subfigura-2}
  \end{subfigure}%

  % A linha em branco acima termina o parágrafo e coloca as próximas
  % sub-figuras em uma nova linha.
  \bigskip % Adiciona um espaço vertical
  % As próximas sub-figuras terão 0.3\textwidth de largura cada,
  % ocupando 90% da largura do texto. Os outros 10% serão preenchidos
  % por espaços em branco de mesmo tamanho (\hfill)
  \begin{subfigure}{0.3\textwidth}
    \includegraphics[width=\textwidth]{example-image}
    \caption{Legenda da sub-figura 3\strut}
    \label{label-da-subfigura-3}
  \end{subfigure}% <- '%' para eliminar o espaço
  \hfill % o \hfill preenche igualmente o espaço em branco em uma linha
  \begin{subfigure}{0.3\textwidth}
    \includegraphics[width=\textwidth]{example-image}
    \caption{Legenda da sub-figura 4\strut}
    \label{label-da-subfigura-4}
  \end{subfigure}%
  \hfill
  \begin{subfigure}{0.3\textwidth}
    \includegraphics[width=\textwidth]{example-image}
    \caption{Legenda da sub-figura 5\strut}
    \label{label-da-subfigura-5}
  \end{subfigure}

  \caption{Legenda da figura inteira}
  \label{label-da-figura-inteira}
\end{figure}

\section{Tabelas}

As tabelas, assim como as figuras, devem aparecer o mais próximo
possível da sua referência no texto. Estas, assim como sua legenda,
devem ser centralizadas. Entretanto, ao contrário das figuras, a
legenda deve aparecer acima da tabela. Um exemplo de tabela é
apresentado na tabela~\ref{tab:exemplo}. Tabelas, assim como figuras,
devem ser referenciadas no texto como tabela~\ref{tab:exemplo} ou
entre parênteses, como (tabela~\ref{tab:exemplo}).

\begin{table}
  \centering
  \caption{Desempenho dos dos diferentes métodos de planejamento da FFTW
    e estimativa do tempo total para a simulação de estudo de caso.}
  \label{tab:exemplo}
  \begin{tabular}{r c c c}
    \toprule
      & Planejamento (\si{\second})
        & Execução (\si{\second})
          & Estimativa (\si{\hour}) \\
    \midrule
      \verb|FFTW_ESTIMATE|   & $\approx0$   & \num{0,598} & \num{5,742} \\
      \verb|FFTW_MEASURE|    & \num{0,190}  & \num{0,508} & \num{4,944} \\
      \verb|FFTW_PATIENT|    & \num{3,276}  & \num{0,433} & \num{4,214} \\
      \verb|FFTW_EXHAUSTIVE| & \num{40,909} & \num{0,433} & \num{4,214} \\
    \bottomrule
  \end{tabular}
\end{table}

Em tabelas deve-se evitar excesso de linhas horizontais, e não se deve
utilizar linhas verticais. Linhas horizontais devem ser usadas para
delimitar a tabela acima e abaixo, e delimitar o cabeçalho e seções
lógicas da tabelas. Sequencias de linhas com a mesma informação não
devem ser separadas por linhas (ver exemplo da
tabela~\ref{tab:exemplo}). Sugere-se utilizar o pacote
\textsf{booktabs}, que define os comandos \verb|\toprule|,
\verb|\midrule|, e \verb|\bottomrule| (entre outros). Consulte o
manual do pacote \textsf{booktabs} para mais informações.

\section*{Agradecimentos}

Agradecimentos são opcionais. Se inseridos, devem estar ao final do
corpo do texto, antes das Referências. Como não fazem parte do corpo
do texto, a seção de agradecimantos não deve ser numerada (use o
estilo Unnumbered Heading).

% Ao final do documento, use \printbibliography para
% imprimir a lista de referências bibliográficas
% (certifique-se de executar o BibTeX antes):
\printbibliography

\end{document}
